\chapter{Conclusiones}

El desarrollo de este proyecto integrador nos permitió contrastar de forma empírica cuatro paradigmas de programación paralela aplicados a un problema real de procesamiento de imágenes, revelando tanto las fortalezas como las limitaciones inherentes a cada modelo, y al enfrentarse a diferentes escalas de cómputo. Los principales aprendizajes se sintetizan a continuación.

\begin{itemize}
    \item \textbf{OpenMP demostró ser la arquitectura más eficiente}, debido a la ausencia total de latencia de comunicación en red, y aprovechando al máximo los 32~núcleos disponibles dentro de un único nodo. Nuestra elección de planificación estática, combinada con la política de \textit{first-touch} para la asignación de memoria en arquitecturas NUMA, permitió maximizar la localidad caché y eliminar prácticamente todo el \textit{overhead} de sincronización entre threads.
    \item \textbf{MPI puro sufrió una penalización altísima por su inherente comunicación}. Con 128~núcleos intercambiando mensajes de forma simultánea, la comunicación entre procesos generado por las operaciones de \texttt{MPI\_Scatterv}, \texttt{MPI\_Gatherv} y el intercambio de halos inhibió el beneficio de la paralelización en la mitad de las configuraciones evaluadas. La alta variabilidad observada en los tiempos de ejecución entre iteraciones (con picos que casi duplicaban el tiempo secuencial) constituyen un factor imposible de predecir y controlar en entornos distribuidos.
    \item \textbf{El modelo Híbrido se posicionó como la solución más equilibrada}, combinando poca comunicación entre procesos (4~procesos MPI) con un máximo aprovechamiento de recursos (32~hilos OpenMP por nodo). Si bien sus valores absolutos fueron modestos comparados con OpenMP, la versión Híbrida superó consistentemente tanto al modelo secuencial como a MPI puro en todas las configuraciones, siendo la alternativa más escalable cuando el problema excede la capacidad de un solo nodo.
    \item Al variar gradualmente la cantidad de hilos de 1 a 32 sobre la configuración más exigente ($5000 \times 5000$, filtro de $5 \times 5$, 9~niveles), \textbf{OpenMP confirmó su dominio como paradigma paralelo}. La versión Híbrida, en cambio, demostró que la comunicación de red entre los 4~nodos MPI es un bottleneck.
    \item Del análisis de eficiencia se obtiene que aumentar la cantidad de núcleos solo se traduce en mejora computacional cuando la fracción paralela del algoritmo domina sobre la fracción de comunicación y sincronización.
\end{itemize}

